% ============================================================================
% 09_conclusion.tex — Conclusion (sec:conclusion), ≤0.5pp per plan §3.
% Merge pass: auctions main + matching-as-robustness phrasing; sharpened
% invariance-content claim; cardinal verdict as an honest boundary; frontier
% migration (not compression) per the genuine gpt-5-mini cells.
% Cut per plan §7: no nation-scale voting, no risk-aversion claim,
% no combinatorial-auction futures.
% ============================================================================

\section{Conclusion}\label{sec:conclusion}

Which levers available to a mechanism designer keep dominant-strategy play intact for boundedly rational agents? Across sealed-bid auctions---the paper's main domain---and four LLM families, with the pattern reproduced in a school-choice matching domain as a robustness check (Appendix~\ref{app:da}), the answer is a stable hierarchy of tiers. Obviously strategy-proof extensive forms sit at the top. A strong second tier contains levers that change no rule of the game, and its content is sharper than ``more explanation helps'': the active ingredient is the \emph{invariance}. Honest text that exposes why truth-telling is safe, and scaffolds that lay out the mechanism's payoff contingencies, deliver large gains---while restatements of the mechanics carry none of it, coming in null on average and sign-mixed in auctions and actively harmful in matching when the invariance statement is omitted. Below the baseline, scaffolds that prompt agents to plan ahead or to model opponents actively backfire, most severely for the weakest models. The ordering parallels the theoretical simplicity hierarchy built for human reasoners, and it localizes the binding constraint: not \emph{computing} the dominant strategy, but \emph{seeing} that truth-telling is safe.

The human experimental record agrees in sign on every rung where it speaks, at the pooled level---clock formats, clock-framing (in three of four families), the menu restatement's null, invariance-carrying versus mechanics-only descriptions, the failing baselines---even though the two populations fail in opposite directions, which is precisely why the agreement is informative: the constraints, and the levers that relieve them, belong to the strategic problem rather than to any one kind of mind. The agreement is ordinal, and we measure that boundary rather than assume it: no single tuning of the agents---temperature, risk personas, framings, explanations---produces cardinal agreement with human play across mechanisms, because level calibration is mechanism- and model-specific and the one knob that restores the human direction of error in second- and third-price auctions breaks first-price play. The testbed ranks levers; it does not simulate levels. Where the human record is silent, the ranking speaks first: its sharpest untested claim is that safety-exposing descriptions will deliver for human subjects the behavioral gains that mechanical restatements have not---a prediction testable in a single incentivized session at a small fraction of the cost of the synthetic program that generated it.

Three directions strike us as most valuable. First, frontier tracking: the constraints migrate rather than vanish as capability grows---the four frontier models we test master the dominant-strategy sealed-bid auction, yet a content-free menu restatement still collapses one of them, and the third-price format still defeats the frontier where we can measure it---so the ranking should be re-estimated on the formats and presentation margins where reasoning still binds as deployed fleets evolve. Second, combining complementary levers---pairing a safety-exposing description with a payoff-tree scaffold---to trace the minimal intervention that closes the gap to the rational benchmark. Third, the method itself: cognitive scaffolds as diagnostic probes, together with the heterogeneity we document across model families, offer a general way to characterize the structure of an artificial agent's reasoning failures---in mechanism design and beyond.
